\chapter{Quantum Hamming Bound}
%%%%%%%%%%%%%%%%%%%%%%%%%%%%%%%%%%%%%%%%%%%%%%%%%%%%%%%%%%%%%%%%%%%%%%%%%%%%%%%
\section{Overview}

This chapter proves the \textbf{quantum Hamming bound} for non-degenerate
qubit stabilizer codes.  Working in the Pauli-string formalism, it shows that
if a quantum code $C \le \mathcal{H}_n$ satisfies the Knill--Laflamme
conditions for $t$-error correction and is non-degenerate, then
\[
\left(\sum_{i=0}^{t}\binom{n}{i}3^i\right)\cdot\dim_{\mathbb{C}}(C) \;\le\; 2^n.
\]

%%%%%%%%%%%%%%%%%%%%%%%%%%%%%%%%%%%%%%%%%%%%%%%%%%%%%%%%%%%%%%%%%%%%%%%%%%%%%%%
\section{Pauli strings}

\begin{definition}[Pauli basis]
\lean{PauliBasis}
\leanok
The four Pauli basis elements $\{I,X,Y,Z\}$, represented as an inductive type.
\end{definition}

\begin{definition}[Pauli string]
\lean{PauliString}
\leanok
A \emph{Pauli string} of length $n$ is a function $p : \mathrm{Fin}\,n \to \mathrm{PauliBasis}$.
\end{definition}

\begin{definition}[Support and weight]
\lean{support}
\lean{weight}
\leanok
The \emph{support} $\mathrm{supp}(p) \subseteq \mathrm{Fin}\,n$ consists of coordinates
where $p(i) \neq I$; the \emph{weight} is $\mathrm{wt}(p) = |\mathrm{supp}(p)|$.
\end{definition}

\begin{definition}[Pauli error set]
\lean{PauliErrorsLe}
\leanok
\uses{PauliString, weight}
The set of all Pauli strings of weight at most $t$:
\[
\mathcal{E}(n,t) = \bigl\{p : \mathrm{PauliString}\,n \;\big|\; \mathrm{wt}(p) \le t\bigr\}.
\]
\end{definition}

\begin{theorem}[Cardinality of the Pauli error set]
\lean{card_pauliErrorsLe}
\leanok
\uses{PauliErrorsLe, weight}
\[
|\mathcal{E}(n,t)| \;=\; \sum_{i=0}^{t}\binom{n}{i}\,3^i.
\]
\end{theorem}

%%%%%%%%%%%%%%%%%%%%%%%%%%%%%%%%%%%%%%%%%%%%%%%%%%%%%%%%%%%%%%%%%%%%%%%%%%%%%%%
\section{$n$-qubit Hilbert space}

\begin{definition}[$n$-qubit Hilbert space]
\lean{Hn}
\leanok
The $n$-qubit Hilbert space
$\mathcal{H}_n = \ell^2\!\bigl(\{0,1\}^n,\mathbb{C}\bigr)$,
implemented as \lean{EuclideanSpace ℂ (Fin n → Fin 2)}.
\end{definition}

\begin{lemma}[Dimension of $\mathcal{H}_n$]
\lean{finrank_Hn}
\leanok
\uses{Hn}
$\dim_{\mathbb{C}}(\mathcal{H}_n) = 2^n$.
\end{lemma}

\begin{definition}[Pauli operator]
\lean{pauliOp}
\leanok
\uses{PauliString, Hn}
For $p \in \mathrm{PauliString}\,n$, the associated Pauli operator
$\hat{p} : \mathcal{H}_n \to \mathcal{H}_n$.
\end{definition}

%%%%%%%%%%%%%%%%%%%%%%%%%%%%%%%%%%%%%%%%%%%%%%%%%%%%%%%%%%%%%%%%%%%%%%%%%%%%%%%
\section{Knill--Laflamme conditions}

\begin{definition}[Knill--Laflamme condition]
\lean{KnillLaflamme}
\leanok
\uses{pauliOp}
A subspace $C \le \mathcal{H}_n$ satisfies the \emph{Knill--Laflamme condition}
for $t$-error correction if for all Pauli strings $E,F$ with $\mathrm{wt}(E),\mathrm{wt}(F) \le t$
there exists $\lambda_{EF}\in\mathbb{C}$ such that
$P_C\,E^\dagger F\,P_C = \lambda_{EF}\,P_C$,
where $P_C$ is the orthogonal projection onto $C$.
\end{definition}

\begin{definition}[Non-degenerate code]
\lean{IsNondegenerate}
\leanok
\uses{KnillLaflamme}
A code is \emph{non-degenerate} if it satisfies the Knill--Laflamme condition
and additionally $P_C E^\dagger F P_C = 0$ whenever $E \neq F$.
\end{definition}

%%%%%%%%%%%%%%%%%%%%%%%%%%%%%%%%%%%%%%%%%%%%%%%%%%%%%%%%%%%%%%%%%%%%%%%%%%%%%%%
\section{Error sphere and sphere-packing}

\begin{definition}[Error sphere]
\lean{ErrorSphere}
\leanok
\uses{pauliOp}
The \emph{error sphere} $\mathrm{ES}(C,t)$ is the subspace
$\bigvee_{\mathrm{wt}(p)\le t} \hat{p}(C)$, i.e.\ the supremum of the
Pauli images of $C$ over all $t$-errors.
\end{definition}

\begin{lemma}[Error subspaces are pairwise orthogonal]
\lean{error_subspaces_orthogonal}
\leanok
\uses{IsNondegenerate, pauliOp}
If $C$ is non-degenerate and $E \neq F$ both have weight $\le t$, then
$\hat{E}(C) \perp \hat{F}(C)$.
\end{lemma}

\begin{lemma}[Dimension of the error sphere]
\lean{error_sphere_dimension}
\leanok
\uses{ErrorSphere, error_subspaces_orthogonal, card_pauliErrorsLe}
If $C$ is non-degenerate,
\[
\dim(\mathrm{ES}(C,t)) \;=\; |\mathcal{E}(n,t)|\cdot\dim(C).
\]
\end{lemma}

%%%%%%%%%%%%%%%%%%%%%%%%%%%%%%%%%%%%%%%%%%%%%%%%%%%%%%%%%%%%%%%%%%%%%%%%%%%%%%%
\section{Quantum Hamming bound}

\begin{theorem}[Quantum Hamming bound]
\lean{quantum_hamming_bound}
\leanok
\uses{error_sphere_dimension, finrank_Hn}
If $C \le \mathcal{H}_n$ is non-degenerate, then
\[
\left(\sum_{i=0}^{t}\binom{n}{i}\,3^i\right)\cdot\dim_{\mathbb{C}}(C)
\;\le\; 2^n.
\]
\end{theorem}

%%%%%%%%%%%%%%%%%%%%%%%%%%%%%%%%%%%%%%%%%%%%%%%%%%%%%%%%%%%%%%%%%%%%%%%%%%%%%%%
%%% added by the blueprint pipeline
\section{Additional declarations}

\begin{definition}[Pauli $I$ matrix]
\lean{sigmaI}
\leanok
The $2\times 2$ identity matrix over $\mathbb{C}$, i.e.\ the Pauli operator $I$.
\end{definition}

\begin{definition}[Pauli $X$ matrix]
\lean{sigmaX}
\leanok
The bit-flip Pauli matrix
$X = \begin{pmatrix} 0 & 1 \\ 1 & 0 \end{pmatrix}$
over $\mathbb{C}$.
\end{definition}

\begin{definition}[Pauli $Y$ matrix]
\lean{sigmaY}
\leanok
The bit-phase-flip Pauli matrix
$Y = \begin{pmatrix} 0 & -i \\ i & 0 \end{pmatrix}$
over $\mathbb{C}$.
\end{definition}

\begin{definition}[Pauli $Z$ matrix]
\lean{sigmaZ}
\leanok
The phase-flip Pauli matrix
$Z = \begin{pmatrix} 1 & 0 \\ 0 & -1 \end{pmatrix}$
over $\mathbb{C}$.
\end{definition}

\begin{definition}[Matrix of a Pauli basis element]
\lean{PauliBasis.toMatrix}
\leanok
\uses{sigmaI, sigmaX, sigmaY, sigmaZ}
Maps each Pauli basis element $\{I,X,Y,Z\}$ to its corresponding $2\times 2$ matrix
$\{\sigma_I,\sigma_X,\sigma_Y,\sigma_Z\}$.
\end{definition}

\begin{definition}[Matrix of a Pauli string]
\lean{pauliMatrix}
\leanok
\uses{PauliBasis.toMatrix, PauliString}
For a Pauli string $p$ on $n$ qubits, the $2^n\times 2^n$ matrix obtained as the
tensor product of the single-qubit Pauli matrices $\sigma_{p(i)}$ over all coordinates~$i$.
\end{definition}

\begin{definition}[Embedding of a non-identity Pauli]
\lean{PauliNZ.toBasis}
\leanok
Embeds a non-identity Pauli ($X$, $Y$, or $Z$) from the three-element type
$\mathrm{PauliNZ}$ into the full Pauli basis $\{I,X,Y,Z\}$.
\end{definition}

\begin{lemma}[Embedded Paulis are non-identity]
\lean{PauliNZ.toBasis_ne_I}
\leanok
\uses{PauliNZ.toBasis}
For every $a \in \mathrm{PauliNZ}$, its image $a.\mathrm{toBasis}$ in the Pauli basis is not the
identity element $I$.
\end{lemma}

\begin{definition}[Pauli string with prescribed support]
\lean{mkWithSupport}
\leanok
\uses{PauliNZ.toBasis, PauliString}
Given a finset $S \subseteq \mathrm{Fin}\,n$ and an assignment $f : S \to \mathrm{PauliNZ}$,
constructs the Pauli string that equals $f(i).\mathrm{toBasis}$ on each $i \in S$ and equals $I$
elsewhere.
\end{definition}

\begin{lemma}[Support of \texttt{mkWithSupport}]
\lean{support_mkWithSupport}
\leanok
\uses{PauliNZ.toBasis_ne_I, mkWithSupport, support}
The Pauli string $\mathrm{mkWithSupport}\,S\,f$ has support exactly $S$.
\end{lemma}

\begin{definition}[Pauli strings with exact support]
\lean{pauliStringsExactSupport}
\leanok
\uses{PauliString, support}
The finset of all Pauli strings on $n$ qubits whose support equals exactly the given
finset $S$.
\end{definition}

\begin{lemma}[Count of Pauli strings with exact support]
\lean{card_pauliStringsExactSupport}
\leanok
\uses{PauliNZ.toBasis, mkWithSupport, pauliStringsExactSupport, support, support_mkWithSupport}
The number of Pauli strings with support exactly $S$ is $3^{|S|}$, since each coordinate in $S$
is assigned one of the three non-identity Paulis.
\end{lemma}

\begin{definition}[Quantum code]
\lean{Code}
\leanok
\uses{Hn}
An $n$-qubit quantum code is a subspace of the $n$-qubit Hilbert space $\mathcal{H}_n$.
\end{definition}

\begin{definition}[Adjoint of a Pauli operator]
\lean{pauliOpAdjoint}
\leanok
\uses{Hn, PauliString, pauliMatrix}
The adjoint (Hermitian conjugate) $\hat{p}^\dagger : \mathcal{H}_n \to \mathcal{H}_n$ of the
Pauli operator associated with a Pauli string $p$.
\end{definition}

\begin{definition}[Projection onto a code]
\lean{codeProj}
\leanok
\uses{Hn}
The orthogonal projection of $\mathcal{H}_n$ onto a subspace $C$, viewed as an endomorphism
$P_C : \mathcal{H}_n \to \mathcal{H}_n$.
\end{definition}

\begin{lemma}[Action of the code projection]
\lean{codeProj_apply}
\leanok
\uses{Hn, codeProj}
For any $x \in \mathcal{H}_n$, the value $P_C\,x$ coincides with the orthogonal projection
$\Pi_C(x)$ of $x$ onto $C$.
\end{lemma}

\begin{lemma}[Image of the code projection lies in $C$]
\lean{codeProj_mem}
\leanok
\uses{Hn, codeProj, codeProj_apply}
For any $x \in \mathcal{H}_n$, the projected vector $P_C\,x$ belongs to the code subspace $C$.
\end{lemma}

\begin{lemma}[Projection fixes code vectors]
\lean{codeProj_eq_self_of_mem}
\leanok
\uses{Hn, codeProj}
If $x \in C$, then $P_C\,x = x$.
\end{lemma}

\begin{lemma}[Idempotence of the code projection]
\lean{codeProj_idempotent}
\leanok
\uses{Hn, codeProj, codeProj_eq_self_of_mem, codeProj_mem}
The code projection is idempotent: $P_C(P_C\,x) = P_C\,x$ for all $x \in \mathcal{H}_n$.
\end{lemma}

\begin{theorem}[Quantum Hamming bound, raw form]
\lean{quantum_hamming_bound_raw}
\leanok
\uses{ErrorSphere, Hn, IsNondegenerate, card_pauliErrorsLe, error_sphere_dimension, finrank_Hn}
For a non-degenerate $[[n,k]]$ quantum code $C \le \mathcal{H}_n$ correcting $t$ errors, with
$\dim_{\mathbb{C}}(C) = 2^k$, the dimension of the error sphere does not exceed the ambient
dimension:
\[
\left(\sum_{j=0}^{t}\binom{n}{j}\,3^j\right)\cdot 2^k \;\le\; 2^n.
\]
\end{theorem}
